% References
\begin{thebibliography}{99}
\bibitem{swindle2012} Swindle, M. M., Makin, A., Herron, A. J., Clubb, F. J. \& Frazier, K. S. Swine as models in biomedical research and toxicology testing. \textit{Toxicologic Pathology} \textbf{40}, 325--357 (2012).

\bibitem{lunney2021} Lunney, J. K. \textit{et al.} Importance of the pig as a human biomedical model. \textit{Science Translational Medicine} \textbf{13}, eabd5758 (2021).

\bibitem{bendus2006} Baxter Bendus, A. E., Mayer, J. F., Shipley, S. K. \& Catherino, W. H. Interobserver and intraobserver variation in day 3 embryo grading. \textit{Fertility and Sterility} \textbf{86}, 1608--1615 (2006).

\bibitem{mogelt2020} M{\o}geltoft Kamstrup, K., Markussen, B., Hay-Schmidt, A., Thorup, F. \& Dybdahl Thomsen, P. Staging of porcine embryos: Comparison of Standard Event System-based statistical clusters with a Carnegie-based staging system. \textit{Developmental Dynamics} \textbf{249}, 1259--1273 (2020).

\bibitem{orrahilly2010} O'Rahilly, R. \& M{\"u}ller, F. Developmental stages in human embryos: Revised and new measurements. \textit{Cells Tissues Organs} \textbf{192}, 73--84 (2010).

\bibitem{schachtschneider2021} Schachtschneider, K. M. \textit{et al.} Epigenetic clock and DNA methylation analysis of porcine models of aging and obesity. \textit{GeroScience} \textbf{43}, 2467--2483 (2021).

\bibitem{zhengh2023} Zheng, H., Vijg, J., Fard, A. T. \& Mar, J. C. Measuring cell-to-cell expression variability in single-cell RNA-sequencing data: A comparative analysis and applications to B cell aging. \textit{Genome Biology} \textbf{24}, 238 (2023).

\bibitem{jin2021} Jin, L. \textit{et al.} A pig BodyMap transcriptome reveals diverse tissue physiologies and evolutionary dynamics of transcription. \textit{Nature Communications} \textbf{12}, 3715 (2021).

\bibitem{teng2024} Teng, J. \textit{et al.} A compendium of genetic regulatory effects across pig tissues. \textit{Nature Genetics} \textbf{56}, 112--123 (2024).

\bibitem{simpson2024} Simpson, L. \textit{et al.} A single-cell atlas of pig gastrulation as a resource for comparative embryology. \textit{Nature Communications} \textbf{15}, 5210 (2024).
\end{thebibliography}